\section{Introduction}\label{sec:introduction}
An integration algorithm may be interrupted when its budget is exhausted,
when another simulation finishes, or when an error diagnostic indicates
sufficient accuracy. There is no reason for that stopping point to coincide
with a power of two. This motivates a construction problem for entire
sequences: choose one generator whose successive prefixes have small
integration errors over a range of sample sizes, rather than optimize only
a distinguished terminal point set.

For an integrand $f:[0,1)^d\to\R$, write
\[
 I(f)=\int_{[0,1)^d} f(\vx)\dif\vx,
 \qquad Q_n(f)=\frac1n\sum_{i=0}^{n-1}f(\vx_i).
\]
We index points from zero, consistently with the special-session slides.
A lattice sequence obtained from a radical inverse has particularly useful
structure when $n$ is a power of the base. Established embedded lattice
constructions exploit this structure to make a single generator useful at
multiple preferred sizes~\cite{CKN06,DKP22}. Our focus is the quality of
\emph{all} prefixes up to a finite construction horizon $N$.

Kronecker sequences provide a complementary family. Their points have the
form $\vx_i=\fracpart{i\vz+\vDelta}$, where fractional parts are taken
componentwise. They are naturally defined at every sample size. Their
quality nevertheless depends on the choice of $\vz$, and a generator that
works well in one projection need not work well in many dimensions.
Classical algebraic constructions and recent geometric investigations provide
useful comparators~\cite{Richtmyer51,DGLPS25}. Here we use a common
kernel-discrepancy criterion to formulate a CBC search for both families.

The criterion is the \emph{weighted sum of expected squared discrepancies}
(WSESD). Its expectation is over a uniform random shift; its sum is over
prefix lengths. The coordinate weights of the kernel encode the relative
importance of variables, while the sample-size weights encode which prefix
lengths the construction should favor. These are different modeling choices.
The special-session talk used sample-size weight $n$ and coordinate weights
$\gamma_j=j^{-2}$.

A further development concerns the order in which a lattice sequence is
traversed. Pride's recent notes~\cite{pride-notes} introduce the Newo
scrambling, whose digit permutations depend on less significant input digits.
For a fixed lattice generator, this changes individual prefixes while
preserving their scalar difference multisets. We show directly in the
permutation notation that the expected squared discrepancy at every prefix
is unchanged. This explains why the same CBC generator can be used for
ordinary van der Corput ordering, Gray-code ordering, or a Newo ordering.
The preservation is stronger than equality of their WSESD values alone.

\subsection{Arbitrary prefixes and missing points}
A prefix objective addresses an unplanned stopping point. Arbitrary missing
observations are a different problem and should not be conflated with it.
There is nevertheless a simple stability bound that explains part of the
motivation of the talk. Let $D(P)$ be the worst-case integration error over
the unit ball of a fixed normed function space, with equal weights on a
finite point collection $P$. If $P=A\cup B$ is a disjoint partition with
$|P|=n$ and $|B|=r<n$, then the error functionals satisfy
\[
 (n-r)(I-Q_A)=n(I-Q_P)-r(I-Q_B).
\]
Consequently, if $D(B)\le C_0$ uniformly over the removed collection,
\begin{equation}\label{eq:removal}
 D(A)\le\frac{nD(P)+rC_0}{n-r}.
\end{equation}
For bounded $r$ and $D(P)=O(n^{-p})$, this gives
$D(A)=O(n^{-\min(p,1)})$. It neither proves that a WSESD construction
is optimal after deletion nor justifies intentionally dropping the first
point of a low-discrepancy sequence. Weighted compound rules can address a
different, higher-order arbitrary-size problem~\cite{HKKN12}; this paper
retains equal integration weights.


\par\smallskip
\noindent\codex{[Scope and theory: Decide which quantitative guarantees the
paper will establish for the CBC constructions. Minimizing the aggregate
criterion does not prove a discrepancy rate at every sample size. Align the
introduction's contribution claims with the final theorem statements and
keep arbitrary prefixes distinct from arbitrary missing observations.]}
\par\smallskip
